\uuid{4mf4}
\titre{Matrice/Transposée/Base/Dimension/Espaces supplémentaires}
\auteur{hanani ; mzari}
\organisation{exo7}
\licence{CC BY-NC-SA 4.0}
\datecreate{2020-01-20}
\niveau{L1}
\module{Algèbre}
\chapitre{Matrice}
\sousChapitre{Propriétés élémentaires, généralités}
\difficulte{3}
\theme{Calcul matriciel, Matrice, Transposée, Base, Dimension, Espaces supplémentaires}
\exo7qcmid{lille-2.matrices.019}

\contenu{
\begin{question}
\qclassification{108.01}
\qtags{motcle=Matrice/Transposée/Base/Dimension/Espaces supplémentaires}
On note $ M_2(\Rr)$ l'espace des matrices carrées d'ordre $2$ à coefficients réels. Soit
$$E= \{M \in M_2(\Rr)\mid ^tM = M\} \quad \mbox{et}\quad F= \{M \in M_2(\Rr) \; ; \;  ^tM = -M\},$$
où $^tM$  désigne la transposée de $M$. Quelles sont les assertions vraies ?
\begin{answers}
\good{$E$ est un espace vectoriel de dimension $3$}.
\bad{$E$ est un espace vectoriel de dimension $2$}.
\good{$F$ est un espace vectoriel de dimension $1$}.
\good{$E$ et $F$ sont supplémentaires dans $ M_2(\Rr)$}.
\end{answers}
\begin{explanations} Soit $M=\left(\begin{array}{rc}
a&b\\ c&d\\ \end{array}\right)$ telle que $^tM = M$, alors $b=c$. Donc la famille
$$\left\{\left(\begin{array}{rc}
1&0\\0&0\\
\end{array}\right),\left(\begin{array}{rc}
0&1\\1&0\\
\end{array}\right),\left(\begin{array}{rc}
0&0\\0&1\\ \end{array}\right)\right\}$$
forme une base de $E$. D'où $\dim E = 3$.
\vskip0mm
Soit $M=\left(\begin{array}{rc}a&b\\c&d\end{array}\right)$ telle que $ ^tM = -M$, alors $a=d=0$ et $c=-b$. Donc $\left\{\left(\begin{array}{rc}
0&1\\-1&0\\ \end{array}\right)\right\}$ est une base de $F$ et donc $\dim F=1$.
\vskip0mm
On vérifie que $E\cap F=\{0_E\}$ et en utilisant le théorème de la dimension d'une somme, on déduit que $E$ et $F$ sont supplémentaires dans $ M_2(\Rr)$.
\end{explanations}
\end{question}
}
